% !TEX program = xelatex
\documentclass[11pt,a4paper]{ctexart}

\usepackage[margin=2.5cm]{geometry}
\usepackage{amsmath}
\usepackage{booktabs}
\usepackage{array}
\usepackage{enumitem}
\usepackage{xcolor}
\usepackage{titlesec}
\usepackage{hyperref}
\hypersetup{colorlinks=true,linkcolor=black,urlcolor=blue}
\setlength{\parskip}{0.5em}
\setlength{\emergencystretch}{3em}
\titleformat{\subsection}{\normalsize\bfseries}{}{0em}{}
\titlespacing{\subsection}{0pt}{0.9em}{0.3em}

\title{\bfseries AI 炒股，到底能赚谁的钱？\\\large 一份研究综述讲给普通投资者听的版本}
\author{\large 叨客宁宁}
\date{}

\begin{document}
\maketitle

\noindent 你可能已经听过太多关于“人工智能炒股”的说法：什么深度学习预测涨跌、大模型读新闻选股、强化学习自动交易。它们听上去都很厉害，也都不太像能听懂的样子。这篇东西不打算教你任何技术，只想把一份长达六十多页的学术综述里，真正和研究结论有关的部分，翻译成日常语言。

先说三个可能和你直觉相反的结论。

第一，AI 确实让交易这件事变强了，但它没有改变交易的利润总量。它更像是一台效率更高的机器，而不是一台印钞机。第二，决定你能不能赚到钱的，往往不是模型聪不聪明，而是手续费和买卖时的价格冲击——这笔账经常大到把利润整个吃掉。第三，你在回测报告里看到的漂亮曲线，绝大多数不是骗人，而是有系统性偏差；偏差不是靠“模型再大一点”能修的。

下面逐条讲。

\subsection*{一、市场里的钱，为什么到不了你的口袋}

先接受一个朴素的事实：在市场上赚钱，本质上是靠信息差。你知道一些别人还不知道的东西，在别人反应之前买入，等他反应过来，价格涨上去，你就赚了。

而信息这个东西有个要命的性质：它用一点少一点。一个消息越多人知道、越多人去交易，它带来的价差就越薄。这不是道德问题，是结构问题。学术上有一整套模型讲这件事，其中一个结论很有意思：市场不可能把自己变得“完全有效”，因为如果价格已经反映了一切信息，就没人再有动力花钱去研究信息了；可一旦没人研究，价格又不可能反映一切。所以市场会稳定地停在一种状态——总是留着一部分没被抹平的价差，等着那些研究得最快、最便宜的人去拿。

这意味着两件事。第一，你面对的永远是有利可图的缝隙，所以“AI 炒股能不能赚钱”这个问题本身就问错了；正确的问题是：你相比别人，在获取和使用信息这件事上，是不是更便宜、更快或更强。第二，任何被广泛使用的信号，它的超额收益都会自动变薄。这解释了一个长期现象：学术论文上被验证过的“赚钱方法”，一旦发表、被大家照做，收益就会衰减。

所以当有人跟你说“某个 AI 模型能稳定赚钱”，你要问的第一个问题不是模型多先进，而是：这个钱是从谁手里赚的，为什么这个人到现在还没学会？

\subsection*{二、AI 真正改变的，是“转化效率”}

把交易拆成三段：第一段是把各种杂乱的信息——价格、财报、新闻、订单流水——整理成有用的线索；第二段是根据线索预测接下来可能怎么走；第三段是决定买多少、什么时候买、怎么买、什么时候撤。

AI 在这三段上都能帮上忙，而且帮得不少。它能读懂几千页财报电话会记录，能从一张卫星照片里数出停车场的车，能盯着每秒上万笔的订单流找出规律——这些换成人来做，要么做不动，要么太贵。所以 AI 的价值是真实的，它让“把信息变成仓位”这个过程的效率大幅提升。

但请注意，它提升的是效率，不是总量。整条链路上能榨出来的利润，是由“还有多少信息没被别人用掉”和“交易要付多少摩擦成本”这两头夹住的。AI 再强，也改不了这两头。

一句话：AI 不是印钞机，它是一台更高效的转化装置。装置再好，也得看喂进去的原料还剩多少，以及转化出来的东西够不够付摩擦的账。

\subsection*{三、真正吃掉利润的，往往是你看不见的成本}

这一段最重要，也最容易被忽略。

假设有个信号，每天能带来大约万分之五（0.05\%）的毛收益。听起来很少，但一年下来也有十几个百分点，很不错。问题来了：如果你每天都调仓，买卖一次的固定费用加价格冲击，成本大约要万分之一到万分之几。把每天的换手率乘上成本，你会发现——成本很可能比万分之五还高。于是这个信号在“每天调仓”的节奏下是亏钱的。

如果你把同样的信号改成每周调一次、每月调一次，成本被摊薄，净收益反而转正了。同一个信号，同样的预测能力，仅仅因为交易频率不同，就从亏钱变成赚钱。

这就是为什么“冲击成本”对个人和大资金的意义完全不同。你卖一套房子，挂个价慢慢卖可以卖到市场价；你要一天之内卖出变现，就得打折，量大折得更狠。股票也是一样：你下单的量占当天成交量越大，你自己就把价格往下砸得越多。大资金的这个“打折”损失，往往就是它追不上小资金灵活优势的原因，也是规模做大的天敌。

所以看到任何一份“策略很好”的报告，先翻到成本那一页。如果它没有成本那一页，那这份报告基本可以放一放。

\subsection*{四、回测报告漂亮的三个陷阱}

第一，幸存者偏差和“挑出来的幸运儿”。如果你试了一万种策略，总有几种在过去几年表现特别好。这跟让一万个人扔硬币、总有人连续扔出十次正面是一个道理——赢家不是因为本事，而是因为样本足够多。学术上称之为多重检验问题：你试的次数越多，“看起来有效”就越不可信。可惜的是，绝大多数对外宣传的回测，都没有说明自己试了多少次。

第二，提前看了答案。这是最隐蔽的一类错误。比如一个模型在计算今天的信号时，用了本该明天才知道的财务数据；或者用整段历史的平均值去标准化数据，等于偷偷把未来信息掺了进去。这类错误不会让程序报错，只会让回测变好看。判断的办法只有一个：问清楚“这一步在什么时间点、能拿到哪些信息”。

第三，也是最近两年最严重的一类：大模型可能早就知道答案。如果你用一个 2024 年发布的大模型，去回测 2020 到 2023 年的新闻，而这段新闻恰好在它的训练材料里，那它不是“推理出”了涨跌，而是“背过”了答案。这种作弊不会表现为异常高的准确率，而是表现为一个看起来很合理的高收益，所以特别难被发现。现在很多“大模型选股”的宣传，问题就出在这里。唯一可靠的检验，是只用模型“出生”之后才发生的数据——而这需要时间慢慢积累，急不得。

\subsection*{五、哪些方向更实在，哪些更像故事}

综述把 AI 交易的六个方向都过了一遍，结论并不平均。

相对可信的一类是贴近“交易摩擦”的：比如怎样把一大笔单子拆开慢慢卖、把成本降到最低；比如给买卖双边报价、赚取买卖价差同时管理库存风险。这类工作之所以可信，是因为第一，它们的数据（每一笔订单和报价）相对干净、量足够大；第二，它们的对手是客观的成本，而不是“市场效率”这种摸不着的东西；第三，它们有公开的数据和代码，别人能重跑。

相对可疑的一类是直接预测收益的：尤其是用大模型读新闻然后预测股价这一类。不是说方法一定错，而是它们的结论很难和“数据泄漏”分开，而且几乎没法被独立验证。此外，靠私有数据（比如某些只卖给机构的消费数据、卫星数据）做出来的高收益，外界永远无法确认——因为数据你根本拿不到。这类结果只能算“经验”，不能算“证据”。

还有一类处于中间：用强化学习做交易。它在执行和做市上有效，但在“直接决定买哪些股票”上，证据明显更弱，原因是它极度依赖你给它设计的那个模拟环境——环境建错了，它会比你更彻底地犯错。

\subsection*{六、面对一份“AI 量化”产品的宣传，问这六个问题}

这份综述最后给出了一个很实用的自查清单。当你看到任何一个声称“AI 交易能赚钱”的说法，不管是论文、路演还是产品介绍，都可以拿这六个问题去问：

第一，数据是不是“当时就能看到的”？有没有用到事后才知道的信息。第二，扣掉手续费和冲击成本之后，样本外的结果是多少？第三，一共试了多少种方案？有没有做多重检验的校正。第四，如果用了大模型，能不能证明测试用的数据不是它训练时见过的？第五，跟什么比的？是跟“无脑买入并持有”比，还是跟业内最强的简单方法比。第六，代码和数据口径能不能拿到？

六个问题里只要有一个回答不了，这个结论的可信度就要打折；如果是关键的两三个答不上来，基本可以当作营销材料看待。

\subsection*{七、给普通投资者的几句实在话}

第一，别指望靠 AI 信号自己下场赚超额收益。这份综述里超过一半的篇幅，其实都在讲同一件事：真正难的从来不是找到一个模型，而是在扣掉所有成本、排掉所有偏差之后，剩下的那一点点是否真的存在。全职的专业机构尚且经常做不到，个人更难。

第二，把 AI 当成降低成本的工具，而不是预测涨跌的水晶球。它帮你更快读懂公告、更清楚地看到持仓的风险暴露、更理性地执行既定计划——这些收益是确定的，比“提高预测准确率”实在得多。

第三，凡是承诺收益率的，一律提高警惕。市场的钱来自信息差和风险承担，任何“稳定高收益”的承诺，背后要么是没算成本的回测，要么是没算进去的风险。

第四，如果你确实要评估一个 AI 量化产品，用上面六个问题，并且要求对方给出成本后的、样本外的、能复现的答案。拿不出来，就不必再往下谈。

\subsection*{写在最后}

这份综述的结论，其实不让人兴奋：AI 在金融交易里的确有真实的用处，但它几乎从不以“让你轻松赚钱”的方式出现。它更像电力——电厂发电效率提高了，不代表每个用电的人都能靠卖电发财，但确实所有人都能用上更便宜、更稳的服务。

如果你读完这篇东西，既没有因此高估 AI，也没有因此把它说得一无是处，那我写这些字的目的就达到了。

\vspace{0.6em}
\noindent\rule{\textwidth}{0.4pt}

\noindent\textbf{免责声明}：本文是对一份学术研究综述的通俗转述，仅用于帮助理解研究结论，不构成任何投资建议，也不对任何具体产品或机构作出评价。文中所有数字均为说明性举例。投资有风险，决策请自行判断或咨询具备资质的专业人士。

\end{document}
